Der Ellbogen ist die Drehachse. Sei $d$ der Abstand des Bizeps, $c$ der des Schwerpunkts des Unterarms und $a$ der der Tasche vom Ellbogen. Der Bizeps dreht den Unterarm nach oben, die Gewichtskräfte von Unterarm und Tasche drehen ihn nach unten:
\begin{align}
 F_\mathrm{M}\cdot d &= m_\mathrm{A}\,g\cdot c + m_\mathrm{T}\,g\cdot a
\end{align}
Also
\begin{align}
 F_\mathrm{M} &= \FMF \\
   &= \frac{\mU\cdot\g\cdot\cA+\mT\cdot\g\cdot\aT}{\dM} \\
   &= \FM \approx \result{\FMP{0}{2}}
\end{align}
Das ist etwa das \ratioP{0}{1}-Fache der Gewichtskraft der Tasche: Muskeln wirken an sehr kurzen Hebelarmen.
